$D_n^2=\det(H_n^{\mathsf T}H_n)$, and
\[ (H_n^{\mathsf T}H_n)_{ik}=\sum_{j<2m}\lambda_{i+j}\lambda_{k+j}=c(k-i),\qquad c(t)=\sum_{j\bmod 2m}\lambda_j\lambda_{j+t}, \]
an even function with $c(t+m)=-c(t)$. For $\omega^m=-1$ and $v=(\omega^i)_{i<m}$, $(H_n^{\mathsf T}H_nv)_i=\omega^i\mu(\omega)$ with $\mu(\omega)=\sum_{u<m}c(u)\omega^{-u}=\tfrac12\sum_{u\bmod2m}c(u)\omega^{-u}=\tfrac12|\hat\lambda(\omega)|^2$, $\hat\lambda(\omega)=\sum_{j<2m}\lambda_j\omega^j$, using the $m$-periodicity of $c(u)\omega^{-u}$. The $m$ vectors $v$ for the $m$ roots of $x^m+1$ are independent (Vandermonde), so $\det(H_n^{\mathsf T}H_n)=2^{-m}\prod_\omega|\hat\lambda(\omega)|^2$; with $\omega=\chi(3)$, $\hat\lambda(\omega)=S(\chi)$. This gives Prop.~D(i). An independent check of the normalisation $2^{-m/2}$: the eigenvector identity $H_n^{\mathsf T}H_nv=\tfrac12|S(\chi)|^2v$ was verified in ball arithmetic for $n=2,\dots,7$ with residual $<10^{-143}$, and $\sqrt{\det H_n^{\mathsf T}H_n}$ agrees with $2^{-m/2}\prod|S(\chi)|$ to the precision of the determinant ball (\path{sage/r11_propD_audit.log}).
